% GENERATED FILE. Edit canonical lessons, then run npm run generate:books.
% canonical-source-hash: fnv1a64:1931b113da07904a
% canonical-lessons: SA-C98-atah, SA-C98-tasmat, SA-C98-hetuh, SA-C98-tadartham, SA-C98-yasmat

\chapter{Therefore, So, Cause, For that reason, Because}
\label{ch:sa-therefore-so-cause-for-that-reason-because}
\hlchaptermodality{98}

\begin{chapteropening}
\textbf{By the end of this chapter:} \emph{I can say therefore, so, a cause, for that reason and because.}

Complete the last lesson of chapter 98: I can say therefore, so, a cause, for that reason and because.
\end{chapteropening}

\section[\texorpdfstring{\sk{अतः} (ataḥ)}{अतः (ataḥ)}]{\sk{अतः} (ataḥ) — therefore}

\label{lesson:SA-C98-atah}



\begin{quote}
Before the new one: say the Sanskrit for a reason, then the Sanskrit for because.
\end{quote}

\subsection*{You'll want to know: \sk{अतः}}
\textbf{\sk{अतः}} — \emph{ataḥ} — \textquotedblleft{}therefore\textquotedblright{}.

\begin{cousinweb}
\emph{ataḥ} is built from the pronoun base \emph{a-} \textquotedblleft{}this\textquotedblright{} and the ending \emph{-tas} \textquotedblleft{}from\textquotedblright{}: \textquotedblleft{}from this\textquotedblright{}, hence \textquotedblleft{}therefore\textquotedblright{} (MW). It is the same \emph{-tas} that you met in \emph{yataḥ}.
\end{cousinweb}

\begin{grammarlens}[title={What you've built}]
One more word for this chapter.
\end{grammarlens}

\subsection*{Your turn}
\begin{itemize}
\raggedright
  \item \emph{Say it:} \emph{ataḥ}
  \item \emph{Say it:} \emph{ataḥ}, once more
  \item \emph{Say it:} say \emph{ataḥ}
  \item \emph{Recall:} say the Sanskrit for a reason, then the Sanskrit for because, then say \emph{ataḥ} again
\end{itemize}

\begin{tcolorbox}[breakable,colback=teal!4,colframe=teal!35!black,title={Before you move on}]
What does \sk{अतः} mean? (\textquotedblleft{}Therefore\textquotedblright{}.) Say it once more.
\end{tcolorbox}
\section[\texorpdfstring{\sk{तस्मात्} (tasmāt)}{तस्मात् (tasmāt)}]{\sk{तस्मात्} (tasmāt) — so, from that}

\label{lesson:SA-C98-tasmat}



\begin{quote}
Before the new one: say the Sanskrit for because, then the Sanskrit for therefore.
\end{quote}

\subsection*{You'll want to know: \sk{तस्मात्}}
\textbf{\sk{तस्मात्}} — \emph{tasmāt} — \textquotedblleft{}so, from that\textquotedblright{}.

\begin{cousinweb}
\emph{tasmāt} is the ablative, \textquotedblleft{}from that\textquotedblright{}, of the pronoun \emph{tad} \textquotedblleft{}that\textquotedblright{} (MW), so \textquotedblleft{}on that account, therefore\textquotedblright{}. The pronoun is inherited from PIE \emph{*tód} (EtymDB), the same word as English \emph{that}.
\end{cousinweb}

\begin{grammarlens}[title={What you've built}]
One more word for this chapter.
\end{grammarlens}

\subsection*{Your turn}
\begin{itemize}
\raggedright
  \item \emph{Say it:} \emph{tasmāt}
  \item \emph{Say it:} \emph{tasmāt}, once more
  \item \emph{Say it:} say \emph{tasmāt}
  \item \emph{Recall:} say the Sanskrit for because, then the Sanskrit for therefore, then say \emph{tasmāt} again
\end{itemize}

\begin{tcolorbox}[breakable,colback=teal!4,colframe=teal!35!black,title={Before you move on}]
What does \sk{तस्मात्} mean? (\textquotedblleft{}So, from that\textquotedblright{}.) Say it once more.
\end{tcolorbox}
\section[\texorpdfstring{\sk{हेतुः} (hetuḥ)}{हेतुः (hetuḥ)}]{\sk{हेतुः} (hetuḥ) — a cause}

\label{lesson:SA-C98-hetuh}



\begin{quote}
Before the new one: say the Sanskrit for therefore, then the Sanskrit for so.
\end{quote}

\subsection*{You'll want to know: \sk{हेतुः}}
\textbf{\sk{हेतुः}} — \emph{hetuḥ} — \textquotedblleft{}a cause\textquotedblright{}.

\begin{cousinweb}
\emph{hetu} comes from the root √hi \textquotedblleft{}urge on, set in motion\textquotedblright{} (CDIAL 14160): MW's first gloss is \textquotedblleft{}impulse\textquotedblright{}, the push that makes a thing happen, and from that \textquotedblleft{}motive, cause, reason\textquotedblright{}. Pali kept it unchanged as \emph{hetu} (CDIAL 14160).
\end{cousinweb}

\begin{grammarlens}[title={What you've built}]
One more word for this chapter.
\end{grammarlens}

\subsection*{Your turn}
\begin{itemize}
\raggedright
  \item \emph{Say it:} \emph{hetuḥ}
  \item \emph{Say it:} \emph{hetuḥ}, once more
  \item \emph{Say it:} say \emph{hetuḥ}
  \item \emph{Recall:} say the Sanskrit for therefore, then the Sanskrit for so, then say \emph{hetuḥ} again
\end{itemize}

\begin{tcolorbox}[breakable,colback=teal!4,colframe=teal!35!black,title={Before you move on}]
What does \sk{हेतुः} mean? (\textquotedblleft{}A cause\textquotedblright{}.) Say it once more.
\end{tcolorbox}
\section[\texorpdfstring{\sk{तदर्थम्} (tadartham)}{तदर्थम् (tadartham)}]{\sk{तदर्थम्} (tadartham) — for that reason}

\label{lesson:SA-C98-tadartham}



\begin{quote}
Before the new one: say the Sanskrit for so, then the Sanskrit for a cause.
\end{quote}

\subsection*{You'll want to know: \sk{तदर्थम्}}
\textbf{\sk{तदर्थम्}} — \emph{tadartham} — \textquotedblleft{}for that reason\textquotedblright{}.

\begin{cousinweb}
\emph{tadartham} is a compound of \emph{tad} \textquotedblleft{}that\textquotedblright{} and \emph{artha} \textquotedblleft{}aim, purpose, sake\textquotedblright{}, used in the accusative as an adverb: \textquotedblleft{}for that end, on that account\textquotedblright{} (MW). The pronoun \emph{tad} is inherited from PIE \emph{*tód} (EtymDB), the same word as English \emph{that}.
\end{cousinweb}

\begin{grammarlens}[title={What you've built}]
One more word for this chapter.
\end{grammarlens}

\subsection*{Your turn}
\begin{itemize}
\raggedright
  \item \emph{Say it:} \emph{tadartham}
  \item \emph{Say it:} \emph{tadartham}, once more
  \item \emph{Say it:} say \emph{tadartham}
  \item \emph{Recall:} say the Sanskrit for so, then the Sanskrit for a cause, then say \emph{tadartham} again
\end{itemize}

\begin{tcolorbox}[breakable,colback=teal!4,colframe=teal!35!black,title={Before you move on}]
What does \sk{तदर्थम्} mean? (\textquotedblleft{}For that reason\textquotedblright{}.) Say it once more.
\end{tcolorbox}
\section[\texorpdfstring{\sk{यस्मात्} (yasmāt)}{यस्मात् (yasmāt)}]{\sk{यस्मात्} (yasmāt) — because}

\label{lesson:SA-C98-yasmat}



\begin{quote}
Before the new one: say the Sanskrit for a cause, then the Sanskrit for for that reason.
\end{quote}

\subsection*{You'll want to know: \sk{यस्मात्}}
\textbf{\sk{यस्मात्}} — \emph{yasmāt} — \textquotedblleft{}because\textquotedblright{}.

\begin{cousinweb}
\emph{yasmāt} is the ablative, \textquotedblleft{}from which\textquotedblright{}, of the relative pronoun \emph{ya-} \textquotedblleft{}which\textquotedblright{} (MW), so \textquotedblleft{}for which reason, because\textquotedblright{}. It answers \emph{tasmāt} \textquotedblleft{}therefore\textquotedblright{}, and MW compares the relative \emph{ya-} with Greek \emph{hos} \textquotedblleft{}who\textquotedblright{}.
\end{cousinweb}

\begin{grammarlens}[title={What you've built}]
One more word for this chapter.
\end{grammarlens}

\subsection*{Your turn}
\begin{itemize}
\raggedright
  \item \emph{Say it:} \emph{yasmāt}
  \item \emph{Say it:} \emph{yasmāt}, once more
  \item \emph{Say it:} say \emph{yasmāt}
  \item \emph{Recall:} say the Sanskrit for a cause, then the Sanskrit for for that reason, then say \emph{yasmāt} again
\end{itemize}

\begin{tcolorbox}[breakable,colback=teal!4,colframe=teal!35!black,title={Before you move on}]
What does \sk{यस्मात्} mean? (\textquotedblleft{}Because\textquotedblright{}.) Say it once more.
\end{tcolorbox}
